% Generated by paper-covert from output/manuscript/sections_revised/07_discussion.md

\section{Discussion}\label{discussion}

\subsection{A presence-detection failure, not a classification failure}\label{a-presence-detection-failure-not-a-classification-failure}

The central interpretive result of this paper is that the map's dominant disagreement with independent ABS statistics is a presence-detection problem rather than a species-classification problem. Three independent lines of evidence converge on this reading. First, canola \emph{composition} matches ABS closely (r=0.82) while total mapped \emph{area} over-calls ABS by 1.47-1.59x, and dividing out the over-call recovers near-exact canola-share agreement (Section 5.2) --- the kind of pattern a presence problem produces (too much land called ``crop-present,'' with the crop mix among that land roughly right) rather than a classification problem (the right amount of land called crop, with the wrong species assigned to some of it). Second, the same over-call magnitude and the same class concentration (Cereal and Legume, not Canola) reproduce independently at national and 100 km-regional scale, which a spurious artefact of one particular region or one particular year would not do. Third, the independent WorldCereal and NLUM spatial comparisons (Section 5.5), built from entirely separate data and with no connection to the ABS validation, corroborate the same diagnosis from a different angle: presence agreement against WorldCereal is symmetric and moderate (73\% both directions) rather than one-sidedly poor, and NLUM's independent grazing-probability layer shows abstained land is systematically more grazing-like than classified land --- consistent with genuine non-cropped land being the dominant source of the abstained population, rather than the classifier failing to recognize crops it should have recognized.

\subsection{The phenology-gate result as a general caution}\label{the-phenology-gate-result-as-a-general-caution}

The phenology-shape presence gate (Section 5.4) is, on paper, the more principled fix: it requires the shape evidence a real crop's growing season should produce (green-up, then senescence) rather than treating any sufficiently large seasonal amplitude as evidence of cropping. It also worked, on its own target metric --- the regional area ratio fell from 1.59 to 1.00, essentially eliminating the area-inflation problem this paper otherwise reports as its most consequential limitation. We rejected it anyway, because the metric it was evaluated against first (area ratio) was not the metric that mattered most for the map's usefulness (per-crop presence recall), and only checking recall per crop, not pooled, revealed that the fix worked by disproportionately removing real canola paddocks --- the map's best-classified and most ABS-consistent crop. We report this as a transferable methodological caution for presence-gate design in crop mapping generally, independent of this specific pipeline: a fix that is more theoretically principled, and that succeeds on its own headline target metric, can still be the wrong choice to ship, if evaluating it on a single pooled metric conceals which sub-population absorbs the improvement's cost. The appropriate response to a design choice like this is not to assume the more sophisticated option is better because it is more sophisticated, but to evaluate it on the metric that actually matters for the map's use, broken out by the sub-populations a pooled metric can hide.

\subsection{Foundation-model embeddings did not help here}\label{foundation-model-embeddings-did-not-help-here}

The Presto negative result (Section 5.1, Section 2.1) is worth discussing on its own terms rather than only as a methods footnote, because it runs against a common framing in the current crop-mapping literature that foundation-model embeddings straightforwardly outperform hand-engineered features. At this label volume (2,194 field-verified trials) and this class granularity (three broad groups), they did not: Presto alone underperformed hand-built spectral indices, and combining the two gave no measurable improvement over the hand-built features alone. We do not read this as evidence that Presto's architecture is unsuited to crop mapping in general --- three of its nine expected input channel groups (Sentinel-1 radar, ERA5 climate, SRTM terrain) were unavailable in this deployment and had to be passed as masked rather than genuinely observed, a specific, testable, and plausible explanation for the gap that this paper's data availability did not allow us to close. The result is reported as a data point against unconditional ``foundation models win'' claims in this literature, conditioned explicitly on label volume, class granularity, and input-channel completeness, not as a general verdict on foundation-model embeddings for crop mapping.

\subsection{The abstain-reason and confidence design as the map's actual answer to area inflation}\label{the-abstain-reason-and-confidence-design-as-the-maps-actual-answer-to-area-inflation}

Because every polygon in the national map carries either a predicted class or one of four explicit abstain reasons, plus a per-polygon NDVI-amplitude value and classification confidence, a downstream user who needs area-accurate rather than composition-accurate output already has the information required to filter toward a tighter, more conservative view of the map without waiting for a future methodological fix. This design choice was itself informed by an earlier failure: an initial design considered building an explicit absence-labelled ``non-crop'' class from presence-only farmer-recorded grazing records, and abandoned it after finding that 78.5\% of a deliberately hard stratum of grazing paddocks passed the same crop-presence test as known crop paddocks --- the negative labels themselves were not reliably negative, because presence-only farmer records cannot straightforwardly be inverted into confirmed absence labels. The abstain-reason and confidence fields shipped instead are the map's considered response to that failure mode: rather than force a binary crop/non-crop decision the available negative-label data could not support reliably, the map reports its uncertainty explicitly, on every polygon, as a first-class output field rather than a caveat buried in accompanying documentation.

\subsection{A covariate's value is task-dependent, not pipeline-dependent}\label{a-covariates-value-is-task-dependent-not-pipeline-dependent}

Two further findings, obtained after the results above were first drafted, add a second general caution alongside the phenology-gate result (Section 6.2). Sentinel-1 was re-tested against the current best classifier candidate --- the shipped model's three indices plus six band-derived indices from \citet{sharma2026wa}, an internally tested but not-yet-adopted feature set discussed further in Section 7 --- rather than the retired three-index baseline against which it was originally measured (Section 5.4, Section 6.3). Against this stronger baseline, Sentinel-1 regresses classifier accuracy by 0.023 macro F1 (temporal) and 0.018 (spatial), concentrated specifically on the Legume class (F1 0.82 to 0.77) --- the same class Sentinel-1 originally appeared to help. Permutation importance explains the reversal mechanistically: the Sharma et al.~red-edge index already resolves most of the Legume/Cereal confusion Sentinel-1 previously corrected (importance +0.134 without Sentinel-1), and once that confusion is resolved, Sentinel-1's additional feature columns have little signal left to contribute and instead dilute the model on a training set that did not grow to match them (importance collapsing to +0.018, with every other feature family scoring negative). The same evaluation against the model actually shipped for Cereal yield produces the opposite result: Sentinel-1 clearly improves yield prediction, adding 0.092 R\textsuperscript{2} on the yield model's own operational (satellite-plus-year/state, temporal-transfer) bar --- its largest gain in any arm tested. A parallel reversal runs the other way: the same band-derived indices that improve classification regress canola yield prediction when substituted for the shipped model's three simpler indices, with temporal R\textsuperscript{2} falling from 0.206 to 0.012. Table 7 summarizes all three results together.

\begin{longtable}[]{@{}
  >{\raggedright\arraybackslash}p{(\linewidth - 8\tabcolsep) * \real{0.2000}}
  >{\raggedright\arraybackslash}p{(\linewidth - 8\tabcolsep) * \real{0.2000}}
  >{\raggedright\arraybackslash}p{(\linewidth - 8\tabcolsep) * \real{0.2000}}
  >{\raggedright\arraybackslash}p{(\linewidth - 8\tabcolsep) * \real{0.2000}}
  >{\raggedright\arraybackslash}p{(\linewidth - 8\tabcolsep) * \real{0.2000}}@{}}
\toprule\noalign{}
\begin{minipage}[b]{\linewidth}\raggedright
Task
\end{minipage} & \begin{minipage}[b]{\linewidth}\raggedright
Covariate tested
\end{minipage} & \begin{minipage}[b]{\linewidth}\raggedright
Without
\end{minipage} & \begin{minipage}[b]{\linewidth}\raggedright
With
\end{minipage} & \begin{minipage}[b]{\linewidth}\raggedright
Change
\end{minipage} \\
\midrule\noalign{}
\endhead
\bottomrule\noalign{}
\endlastfoot
Classification, macro F1 (temporal) & Sentinel-1, added to the leading classifier candidate & 0.887 & 0.864 & $-$0.023 \\
Classification, macro F1 (spatial) & Sentinel-1, added to the leading classifier candidate & 0.886 & 0.868 & $-$0.018 \\
Cereal yield, R\textsuperscript{2} (temporal, satellite+year/state) & Sentinel-1, added to the shipped yield model & 0.419 & 0.511 & +0.092 \\
Cereal yield, R\textsuperscript{2} (spatial, satellite+year/state) & Sentinel-1, added to the shipped yield model & 0.569 & 0.591 & +0.022 \\
Canola yield, R\textsuperscript{2} (temporal, satellite+year/state) & \citet{sharma2026wa} indices, substituted for the shipped model's three indices & 0.206 & 0.012 & $-$0.194 \\
\end{longtable}


\textbf{Table 7.} Marginal value of Sentinel-1 and the \citet{sharma2026wa} band-derived indices, by downstream task. The same covariate that regresses classification helps Cereal yield; the same substitution that helps classification regresses canola yield. The classification rows' ``without'' baseline (0.887/0.886) is the leading candidate scored on the row-matched subset used for a fair Sentinel-1 A/B comparison (2,133 trials with Sentinel-1 coverage, 497 of 543 test rows), and is therefore not identical to the same candidate's full-population score reported in Table 8 (0.890/0.892, 2,194 trials) --- both are correct for their own comparison, and the \textasciitilde0.003-0.006 difference reflects the smaller, Sentinel-1-restricted row set, not a discrepancy in the underlying model.

We read this as a second general methodological caution, alongside the phenology-gate result (Section 6.2): a covariate or feature set validated as beneficial for one task in a shared feature pipeline should not be assumed beneficial for every task that pipeline feeds, even within the same project and the same underlying satellite inputs. Sentinel-1 backscatter is sensitive to canopy structure and biomass, information a yield regressor needs and a phenology-shape classifier, once its dominant residual confusion is otherwise resolved, does not. The practical implication for this project is that Sentinel-1 acquisition, previously considered primarily as a classifier-accuracy lever, is now better justified --- if pursued at national scale --- as a yield-accuracy lever specifically (Section 7.2).

\subsection{Generalizability across seasons: regional evidence, national open question}\label{generalizability-across-seasons-regional-evidence-national-open-question}

The nine-year 100 km regional test identifies two structurally weaker years --- 2017 and 2018 --- with lower classified share and lower mean NDVI amplitude than the other seven seasons, attributable respectively to Sentinel-2B's mid-2017 operational start (reducing early-mission revisit frequency and cloud-gap robustness) and a documented 2018 drought. This is useful evidence that the pipeline's behaviour is not uniform across seasons, and gives a reason to expect, but not to assume, that a national multi-year extension will show similar season-to-season variation once run. We deliberately do not extrapolate the regional nine-year pattern into a national multi-year claim: a 100 km regional block, however intensively validated, is one climatic and cropping-system context among several across the continent, and whether the same two mechanisms (satellite-mission timing, drought) dominate nationally, or whether other regions show different weak seasons for different reasons, is precisely the open question the planned national multi-year extension (Section 3.2, Section 7) is designed to answer. This paper's own single-national-season (2024) results should accordingly be read as characterizing that one season's performance, not as an implicit claim about multi-year national stability.

\subsection{The Legume classification error remains open}\label{the-legume-classification-error-remains-open}

Distinct from the presence-detection story that dominates this paper's other findings, Legume classification carries its own unresolved error: the national map assigns 18.4\% of classified area to Legume against an ABS share of 8.7\% (Section 5.3), more than double the official figure. This is not an artefact of the map's decision rule --- argmax and mean-predicted-probability agree on essentially the same share (18.4\% and 18.7\% respectively; Section 4.6), which rules out a systematic tie-break or threshold quirk and confirms the classifier genuinely, if wrongly, favours Legume at this rate rather than merely reporting it inconsistently. The confusion-matrix evidence is suggestive but not sufficient on its own: Cereal paddocks are misclassified as Legume more often than the reverse (23 of 279 true Cereal paddocks under temporal transfer; Section 5.1), consistent in direction with the national over-prediction but not obviously large enough in magnitude to fully explain an 8.7-to-18.4-point gap by itself. We report this as the paper's largest genuinely unresolved classification error --- distinct in kind from the area-inflation problem Sections 6.1-6.2 diagnose and address, because it is a species-level confusion (which crop, given that something is present) rather than a presence-level one (whether anything is present at all) --- and leave the mechanism open rather than speculate past what the available diagnostics support (Section 7, Limitation 2).
